%% 05_discussion.tex
% Calibrated to analysis/results.md, analysis/error_taxonomy.md,
% analysis/feature_extraction.md, analysis/recall_at_k.md

\section{Discussion}
\label{sec:discussion}

\paragraph{Interpreting the A $>$ B $>$ C ordering.}
% Source: analysis/results.md, analysis/error_taxonomy.md
System~A outperforms both structured variants on exact-match accuracy by 5.27~pp overall
(vs.\ B) and 7.21~pp overall (vs.\ C), with statistically significant pairwise differences.  The ordering
is real, but its interpretation requires care.  The error taxonomy reveals that
44--58\% of all failures arise from reasoning truncation under the 512-token output
budget, disproportionately affecting B and C because their FHIR-JSON contexts demand
longer reasoning chains before a final answer can be stated.  Had the output budget
been unconstrained, the gap would likely narrow, and would narrow further for C,
whose resource-aware retrieval already demonstrates superior retrieval recall
(0.852 vs.\ 0.762 at $k=5$; Table~\ref{tab:recall}).
The 512-token limit was set to control API costs during the evaluation run and
is a methodological artefact, not a property of narrative vs.\ structured RAG.
Future work should sweep ANSWER\_MAX\_TOKENS to isolate the architecture effect from
the budget effect.

\paragraph{Why the narrative format has a temporal weakness.}
The temporal-anchor failure rate gradient (A:~30\%, B:~18\%, C:~10\%) is the most
structurally informative finding in the error taxonomy.  Narrative chunks do not
contain the question's \texttt{reference\_date} (a per-question parameter that changes
across rows in the question bank), so the LLM must infer a temporal anchor from
surrounding text.  FHIR resources carry explicit date fields
(CarePlan.period.start/end, MedicationAdministration.effectiveDateTime,
MedicationRequest.authoredOn) that enable the LLM to reason about ``when'' even when
the reference date is absent from the prompt context.  This is a structural property
that would persist regardless of token budget, and it motivates future work embedding
reference-date context directly into narrative chunks at indexing time.

\paragraph{The cross-resource family as a diagnostic.}
Structured retrieval should, in theory, have the largest advantage on cross-resource
questions that require traversing FHIR references (e.g.\ linking a CarePlan to its
MedicationRequests to their Medications).  In practice, cross-resource accuracy is
at parity across all three systems (A=43.4\%, B=44.5\%, C=43.4\%;
Table~\ref{tab:family_results}).  Two non-exclusive explanations exist.  First,
the discriminative cross-resource templates surviving the W2 ground-truth patch
(tier label, schedule kind, component presence) are not the most structurally demanding
question type; the trivial templates were removed precisely because they were
non-discriminative.  Second, the reasoning-truncation problem affects all three systems
equally on this family, masking any structural advantage.

\paragraph{The representation dissociation: QA winner $\neq$ ML-feature winner.}
The central finding of the feature-extraction arm (O6) is a \emph{dissociation} between
QA performance and adherence-prediction performance.  System~A wins the QA comparison
(40.6\% vs.\ 33.4--35.3\%), but FS-Structured features substantially outperform
FS-Narrative features on adherence classification (LightGBM AUC 0.997 vs.\ 0.843;
bootstrap CI excludes zero).  FS-Narrative and FS-Aware are tied within the bootstrap
interval ($\Delta{=}+0.060$ [$-0.051$, $+0.177$]).

The practical implication is that the choice of representation should be made
\emph{separately} for each downstream task.  A clinical workflow that uses narrative RAG
for case-manager Q\&A (where A wins) and structured FHIR features for an adherence
risk-stratification model (where FS-Structured wins) would outperform a workflow that
uses a single representation for both.  This finding complements prior work by Nateghi
Haredasht et al.~\citep{nateghi2025predicting} (LLM-extracted narrative features
for MOUD retention) and Gao et al.~\citep{gao2024raw} (raw numeric features
often match or beat LLM embeddings), which together establish that neither representation
uniformly dominates across downstream use cases.

\paragraph{A point where readers may disagree.}
A reasonable reviewer might argue that the near-perfect FS-Structured AUC (0.997--1.000)
renders the O6 comparison trivial: if structured features are essentially tautological
summaries of the labelling rule, the ``dissociation'' finding collapses to ``the features
used to generate the label also predict the label.''  We acknowledge this directly.  The
synthetic-label design (consecutive missed doses OR gap $>$ 1.5$\times$ interval) was
chosen for reproducibility and programmatic verifiability; its alignment with the
FS-Structured feature set is intentional and documented.  The meaningful signal in O6 is
not the absolute FS-Structured AUC, but (1) the tied CI between FS-Narrative and FS-Aware
despite very different feature sources, and (2) the relative ordering (Structured $\gg$
Narrative $\approx$ Aware) which is robust across both classifiers.  We do not claim
that structured features will achieve near-perfect AUC on a real noisy cohort; we claim
that they encode the adherence signal more directly than narrative-derived or
retrieval-trace features, a claim that should hold under richer real-world labels as well.

\paragraph{Practical guidance for system designers.}
For practitioners building RAG over specialty-medication clinical records, the choice of
representation should follow the downstream task.  When the answer LLM's output budget
is constrained, narrative RAG is the lower-risk choice: its compact context (490 mean
tokens-in vs.\ 4,463 for System~B) allows more reasoning headroom within a fixed token
limit.  If output budgets are generous, structured RAG's temporal grounding becomes a
meaningful advantage; the temporal-anchor failure gradient (A:~30\%, B:~18\%, C:~10\%)
suggests that structured systems would recover faster than narrative systems once the
token ceiling is raised.  For downstream adherence risk stratification, structured FHIR
features are preferable to reusing the narrative representation: extracting numeric
adherence metrics directly from FHIR resources is both more interpretable and more
predictive in this evaluation.  Finally, LLM-generated narratives should be treated as
summaries, not as drop-in replacements for the structured record.  The fidelity audit
confirms 100\% entity recovery after prompt iteration, but entity recovery does not
guarantee that every temporal relation is preserved in a form the answer LLM can
exploit during reasoning.
